\section{Occupation bounds and perturbation estimates}
\label{f:local-section}

The next estimates turn the quadratic stability bound into uniform
control of local particle numbers and changes of pressure.  They play
the role of the occupation bounds in the theory of superstable
interactions \citep{Ruelle1970}.  We prove them for the precise family
of deletions and interpolations used in the range comparison.

From this section onward the occupation estimates are used with $t$
fixed.  Their constants may depend on $t$ and on a compact set of
$(\beta,h)$, but never on the number or sizes of the attraction ranges.
The finite volume systems in \eqref{f:modified-energy} are included;
in particular, deleting interactions between larger boxes will not
require a new thermodynamic limit theorem.

If a bounded region meets $k$ small cells, its grand partition
function has the exact decomposition
\begin{equation}\label{f:cell-measure}
 \Xi^v(\beta,h)
 =\sum_{n_1,\ldots,n_k\ge0}
    \frac{\ee^{th\sum_i n_i}}{\prod_i n_i!}
    \int_{\prod_i C_i^{n_i}}\exp(-\beta U^v(X))\,\dd X.
\end{equation}
Indeed, the $N!/\prod_i n_i!$ assignments of $N=\sum_i n_i$ labels
to disjoint cells cancel the global $1/N!$.  This identity also gives
the conditional grand distribution within one cell after the outside
configuration has been specified.

\begin{proposition}[Uniform occupation moments]\label{f:moments}
Fix $t$ satisfying \eqref{f:t-requirements} and a compact
$\mathcal D\subset\Omega$.  There are constants
$C_{\exp},C_1<\infty$, depending only on $t$ and $\mathcal D$, such
that every grand canonical measure for an interaction
\eqref{f:modified-energy} satisfies
\begin{equation}\label{f:moment-estimates}
 \sup_i\mathbb E\exp(\beta D_0n_i^2/4)\le C_{\exp},
 \qquad \sup_{i,j}\mathbb E[n_i n_j]\le C_1
 \qquad ((\beta,h)\in\mathcal D).
\end{equation}
The bounds hold for every bounded region with finitely many cropped
cells, and uniformly for all permitted deletions and interpolations.
\end{proposition}

\begin{proof}
If the region meets no small cell, its grand measure is the vacuum and
the assertion is immediate.  Otherwise, let the number of cells be positive.
Write $\alpha=D_0-M/2$, $u_0=\beta D_0/4$, and enclose
$\mathcal D$ in $\beta\in[\beta_-,\beta_+]$, $\abs{h}\le H$.
We first check, before taking any maximum, that
$F_i:=\mathbb E\exp(u_0n_i^2)$ is finite in each fixed volume.
By \eqref{f:cell-stability}, insertion of this exponential into
\eqref{f:cell-measure} bounds its numerator by the product of $k$
one-variable series with exponents
\[
 -\beta(\alpha-D_0/4)n_i^2
 -\beta\alpha\sum_{j\ne i}n_j^2
 +(th+\beta B_0)\sum_j n_j.
\]
Each spatial integral contributes at most
$\prod_j(b^3)^{n_j}$, with the factors $1/n_j!$ already in
\eqref{f:cell-measure}.  The quadratic coefficients are positive in
magnitude because
\[
 \alpha-D_0/4=3D_0/4-M/2\ge11D_0/16>0.
\]
Thus all these series are finite, and the denominator $\Xi^v$ is at
least one.  In particular,
$F_{\max}:=\max_iF_i<\infty$.  To bound this maximum independently of
the number of cells, we shall prove
\[
 F_{\max}\le C_2 F_{\max}^{\eta_0},
 \qquad \eta_0=4M^2/D_0^2\le1/16<1,
\]
where $C_2$ depends only on $t$ and $\mathcal D$.
The ratio of attraction strength to core repulsion will therefore
make the dependence on neighboring moments sublinear.

Condition on the configuration $Y$ outside cell $i$.  For $n$
particles inside that cell, the part of the energy involving at least
one inside particle is bounded below by
\begin{equation}\label{f:conditional-energy}
 (D_0-J_{ii}/2)n^2-B_0n-nG_i(Y),
 \qquad G_i(Y)=\sum_{j\ne i}J_{ij}n_j.
\end{equation}
This uses the complete core estimate inside the cell, discards
nonnegative core terms across cells, and bounds all attractions by
their cell suprema.  Dividing the conditional weights by the
outside-only Boltzmann factor leaves a denominator containing the
empty-cell term one.  Hence the conditional exponential moment is
at most its unnormalized numerator.  The inequality
\[
 nG_i\le D_0n^2/4+G_i^2/D_0
\]
and the factor $\exp(u_0n^2)$ leave the quadratic coefficient
$D_0/2-J_{ii}/2\ge7D_0/16$.  Consequently
\begin{equation}\label{f:conditional-moment}
 \mathbb E[\exp(u_0n_i^2)\mid Y]
 \le C_2\exp(\beta G_i(Y)^2/D_0),
\end{equation}
where the finite constant
\[
 C_2=\sum_{n\ge0}\frac{b^{3n}}{n!}
       \exp\left(-\frac{7\beta_-D_0}{16}n^2
                    +(tH+\beta_+B_0)n\right)
\]
depends only on $t$ and $\mathcal D$ and is at least one.

To close this estimate, Cauchy--Schwarz with the weights $J_{ij}$
and the row bound give
\[
 G_i^2\le M\sum_{j\ne i}J_{ij}n_j^2.
\]
Set $\theta_j=J_{ij}/M$ for $j\ne i$ and
$\theta_\varnothing=1-\sum_{j\ne i}\theta_j\ge0$.
With the exponent $\eta_0$ chosen above,
$\beta G_i^2/D_0\le\eta_0u_0\sum_{j\ne i}\theta_jn_j^2$.
Apply convexity of the exponential with the additional zero argument
of weight $\theta_\varnothing$, and then the concavity of
$x\mapsto x^{\eta_0}$, to obtain
\[
 \begin{split}
 \mathbb E\exp(\beta G_i^2/D_0)
 &\le\theta_\varnothing+
       \sum_{j\ne i}\theta_j\mathbb E\exp(\eta_0u_0n_j^2)\\
 &\le\theta_\varnothing+\sum_{j\ne i}\theta_jF_j^{\eta_0}
 \le F_{\max}^{\eta_0},
 \end{split}
\]
where the last inequality uses $F_{\max}\ge1$.
Taking expectations in \eqref{f:conditional-moment} and then the
maximum over cells yields
$F_{\max}\le C_2F_{\max}^{\eta_0}$, and therefore
\[
 F_{\max}\le C_2^{1/(1-\eta_0)}=:C_{\exp}.
\]
This bound is independent of the number of cells.  Since
$\exp(u_0n_i^2)\ge u_0n_i^2$ and
$u_0\ge\beta_-D_0/4$, the second moments are uniformly bounded.
Cauchy--Schwarz gives the mixed moment bound in
\eqref{f:moment-estimates}.

Cropping a cell only decreases its volume and the relevant suprema.
Every deletion or interpolation allowed in \eqref{f:modified-energy}
retains the full intra-cell core and the same row bound.  All steps
above therefore apply without change to those systems.
\end{proof}

\begin{corollary}[Absolute perturbation bound]\label{f:perturbation}
In a fixed bounded region, let $v_\theta=(1-\theta)v_0+\theta v_1$,
$0\le\theta\le1$, be bounded symmetric pair interactions satisfying
the hypotheses of \Cref{f:moments} with the same $t$ and $M$.
Let $J^{\mathrm{diff}}$ be a symmetric nonnegative matrix such that
\[
 \abs{v_1(x,y)-v_0(x,y)}\le J^{\mathrm{diff}}_{ij}
 \quad (x\in C_i,\ y\in C_j).
\]
Then, uniformly for $(\beta,h)\in\mathcal D$,
\begin{equation}\label{f:interpolation-bound}
 \abs{\log\Xi^{v_1}(\beta,h)-\log\Xi^{v_0}(\beta,h)}
 \le\frac{\beta C_1}{2}\sum_{i,j}J^{\mathrm{diff}}_{ij}.
\end{equation}
Thus signed deletions are controlled by the absolute interaction
difference, with no cancellation assumption.
\end{corollary}

\begin{proof}
Differentiation in $\theta$ is justified directly in finite volume.
The empty region is immediate.  Put $\alpha=D_0-M/2>0$.
If there are $k\ge1$ cells and the region has volume $V$, the common
energy bound gives
$U^{v_\theta}\ge\alpha N^2/k-B_0N$, while boundedness of the pair
difference gives $\abs{\partial_\theta U^{v_\theta}}\le C N^2$.
Uniformly in $\theta$ and on $\mathcal D$, the differentiated grand
series is therefore dominated by a constant times
\begin{equation}\label{f:derivative-domination}
 \sum_{N\ge0}(1+N^2)\frac{V^N}{N!}
       \exp\left(-\beta_-\alpha N^2/k
                      +(tH+\beta_+B_0)N\right)<\infty.
\end{equation}
The volume and the constant in this domination are fixed only for
the differentiation; the resulting estimate below is uniform in
volume.  Hence
\[
 \frac{\dd}{\dd\theta}\log\Xi^{v_\theta}
       =-\beta\mathbb E_\theta\partial_\theta U^{v_\theta}.
\]
For every configuration,
\[
 \abs{\partial_\theta U^{v_\theta}}
 \le\frac12\sum_{i,j}J^{\mathrm{diff}}_{ij}n_i n_j,
\]
where the diagonal uses $n_i(n_i-1)\le n_i^2$.
Apply \Cref{f:moments} and integrate over $\theta\in[0,1]$.
\end{proof}

\begin{corollary}[Pressure control by attraction mass]
\label{f:tail-pressure}
For a finite or countably infinite member of \eqref{f:family}, let
\[
 \varphi_k=\psi_t-t\sum_{j\le k}\alpha_jw_{R_j},
 \qquad m_k=\sum_{j>k}\alpha_j.
\]
For every compact $\mathcal D\subset\Omega$ there is a constant
$C_{\mathcal D,t}<\infty$, independent of the ranges, their number,
and $k$, such that
\begin{equation}\label{f:tail-pressure-bound}
 0\le Q_\varphi(\beta,h)-Q_{\varphi_k}(\beta,h)
 \le C_{\mathcal D,t}m_k
 \qquad ((\beta,h)\in\mathcal D).
\end{equation}
The same upper bound holds for the difference of normalized finite
volume log grand partition functions in every cube of side $L\ge1$.
\end{corollary}

\begin{proof}
Interpolate the whole tail at once:
\[
 \varphi_{k,\theta}
 =\varphi_k-\theta t\sum_{j>k}\alpha_jw_{R_j},
 \qquad 0\le\theta\le1.
\]
Even for infinitely many terms, the tail converges uniformly and is
a bounded continuous pair interaction, of supremum norm at most
$t\norm{w}_\infty m_k$.  Thus the derivative exists for each finite
configuration and the domination \eqref{f:derivative-domination}
justifies differentiation directly, including the infinite endpoint.
All intermediate attractions obey the original mass bound.

By the proof of \eqref{f:row-bound}, the difference matrix can be
chosen with
\[
 \sup_i\sum_jJ^{\mathrm{diff}}_{ij}\le c_btm_k.
\]
There are at most $216L^3$ cells.  Dividing
\eqref{f:interpolation-bound} by $tL^3$ gives the claimed finite
volume upper bound, for instance with
$C_{\mathcal D,t}=108\beta_+C_1c_b$.  Its lower bound follows because
adding a nonnegative attraction decreases the energy.

\Cref{f:thermodynamics} already supplies thermodynamic pressure
limits for both endpoints, including the infinite sum.  Passing to
those limits proves \eqref{f:tail-pressure-bound}; convergence of
truncation pressures is a consequence, not an input, of this estimate.
\end{proof}
